\chapter{جزئیات کامل الگوریتم تخمین موقعیت و سرعت موانع با فیلتر کالمن}
\label{appendix:kalman_speed}

این پیوست، جزئیات فنی الگوریتم مورد استفاده برای تخمین موقعیت و سرعت نسبی موانع را تشریح می‌کند. در این الگوریتم، ابتدا موانع با استفاده از آشکارساز دوبعدی در تصویر چپ شناسایی می‌شوند. سپس با ترکیب جعبه محدودکننده هر مانع و نقشه اختلاف مکانی دوربین استریو، فاصله و زاویه نسبی مانع نسبت به اسکوتر محاسبه می‌شود. در ادامه، برای کاهش نویز اندازه‌گیری و تخمین سرعت نسبی، از فیلتر کالمن استفاده می‌شود.

ورودی فیلتر کالمن، مختصات دوبعدی مانع در چارچوب لینک پایه اسکوتر است. این مختصات از ترکیب عمق تخمین‌زده‌شده از نقشه اختلاف مکانی و زاویه نسبی مانع نسبت به محور دید دوربین به‌دست می‌آید. خروجی فیلتر شامل موقعیت پالایش‌شده مانع، عمق پالایش‌شده، مؤلفه‌های سرعت نسبی و سرعت کلی نسبی مانع است. باید توجه داشت که سرعت تخمین‌زده‌شده در این روش، سرعت نسبی مانع نسبت به اسکوتر است و نباید به‌عنوان سرعت مطلق مانع در دستگاه مختصات جهانی تفسیر شود.

در پیاده‌سازی انجام‌شده، برای هر مانع معتبر یک فیلتر کالمن مستقل نگه‌داری می‌شود. اگر شناسه ردیابی مانع از خروجی آشکارساز موجود باشد، همان شناسه برای اتصال اندازه‌گیری‌های متوالی استفاده می‌شود. در غیر این صورت، یک کلید جایگزین بر اساس کلاس جسم و موقعیت گسسته‌شده مرکز جعبه محدودکننده ساخته می‌شود. این سازوکار باعث می‌شود هر مانع در فریم‌های متوالی تا حد امکان با وضعیت قبلی خود مرتبط باقی بماند.

%----------------------------------------------------------
\section{بردار حالت}

بردار حالت فیلتر کالمن برای هر مانع به‌صورت چهاربعدی تعریف می‌شود:

\begin{equation}
	\mathbf{x}*k =
	\begin{bmatrix}
		x_k \
		v*{x,k} \
		y_k \
		v_{y,k}
	\end{bmatrix}
	\label{eq:app_kalman_state}
\end{equation}

در این رابطه، ($x_k$) و ($y_k$) مختصات نسبی مانع در صفحه حرکت و نسبت به چارچوب لینک پایه اسکوتر هستند. مؤلفه ($x_k$) راستای طولی و روبه‌روی اسکوتر را نشان می‌دهد و مؤلفه ($y_k$) بیانگر جابه‌جایی جانبی مانع نسبت به محور حرکت اسکوتر است. همچنین ($v_{x,k}$) و ($v_{y,k}$) مؤلفه‌های سرعت نسبی تخمین‌زده‌شده در همان چارچوب هستند.

ترتیب مؤلفه‌ها در بردار حالت مطابق پیاده‌سازی نرم‌افزاری انتخاب شده است؛ به‌طوری‌که مؤلفه سرعت هر راستا بلافاصله پس از مؤلفه مکان همان راستا قرار می‌گیرد. بنابراین، فیلتر کالمن در این کاربرد تنها برای هموارسازی فاصله استفاده نمی‌شود، بلکه موقعیت دوبعدی و سرعت نسبی مانع را به‌صورت هم‌زمان تخمین می‌زند.

%----------------------------------------------------------
\section{مدل حرکت}

برای بازه‌های زمانی کوتاه، حرکت نسبی هر مانع با مدل سرعت ثابت تقریب زده می‌شود. بر این اساس، رابطه پیش‌بینی حالت به‌صورت زیر است:

\begin{equation}
	x_{k+1}=x_k+v_{x,k}\Delta t,
	\quad
	v_{x,k+1}=v_{x,k},
	\qquad
	y_{k+1}=y_k+v_{y,k}\Delta t,
	\quad
	v_{y,k+1}=v_{y,k}.
\end{equation}

بنابراین، ماتریس گذار حالت برابر است با:

\begin{equation}
	\mathbf{F}_k =
	\begin{bmatrix}
		1 & \Delta t & 0 & 0 \\
		0 & 1        & 0 & 0 \\
		0 & 0        & 1 & \Delta t \\
		0 & 0        & 0 & 1
	\end{bmatrix}
	\label{eq:app_kalman_transition}
\end{equation}

در پیاده‌سازی عملی، مقدار ($\Delta t$) برای هر مانع به‌صورت مستقل و بر اساس اختلاف زمان واقعی میان دو به‌روزرسانی متوالی همان مانع محاسبه می‌شود. برای جلوگیری از ناپایداری عددی در شرایطی که فاصله زمانی بسیار کوچک یا بسیار بزرگ باشد، مقدار ($\Delta t$) در بازه زیر محدود می‌شود:

\begin{equation}
	0.01 \leq \Delta t \leq 1.0
	\label{eq:app_kalman_dt_clip}
\end{equation}

این محدودسازی باعث می‌شود فیلتر در برابر تغییرات نرخ پردازش، تأخیرهای موقت و نامنظمی زمان دریافت داده‌ها پایدارتر عمل کند.

%----------------------------------------------------------
\section{مدل اندازه‌گیری}

اندازه‌گیری مستقیم حاصل از زنجیره ادراک شامل موقعیت دوبعدی مانع در چارچوب لینک پایه است. بنابراین بردار اندازه‌گیری به‌صورت زیر تعریف می‌شود:

\begin{equation}
	\mathbf{z}_k =
	\begin{bmatrix}
		x^{meas}_k \
		y^{meas}_k
	\end{bmatrix}
	\label{eq:app_kalman_measurement}
\end{equation}

از آنجا که سرعت مانع به‌صورت مستقیم اندازه‌گیری نمی‌شود، ماتریس اندازه‌گیری فقط مؤلفه‌های مکانی بردار حالت را استخراج می‌کند:

\begin{equation}
	\mathbf{H} =
	\begin{bmatrix}
		1 & 0 & 0 & 0 \
		0 & 0 & 1 & 0
	\end{bmatrix}
	\label{eq:app_kalman_H}
\end{equation}

بنابراین رابطه اندازه‌گیری به‌صورت زیر است:

\begin{equation}
	\mathbf{z}_k = \mathbf{H}\mathbf{x}_k + \mathbf{v}_k
	\label{eq:app_kalman_measurement_model}
\end{equation}

که در آن ($\mathbf{v}_k$) نویز اندازه‌گیری است. این نویز ناشی از عواملی مانند خطای تطبیق استریو، تغییرات کادر تشخیص، نویز تصویر و ناپایداری پیکسل‌های معتبر در نقشه اختلاف مکانی است.

%----------------------------------------------------------
\section{تشکیل اندازه‌گیری از داده استریو}

برای تشکیل اندازه‌گیری، ابتدا جعبه محدودکننده مانع از خروجی آشکارساز دریافت می‌شود. این جعبه روی نقشه اختلاف مکانی متناظر قرار می‌گیرد تا مقادیر اختلاف مکانی مربوط به همان ناحیه استخراج شوند. پیش از نمونه‌برداری از نقشه اختلاف مکانی، ناحیه داخلی جعبه محدودکننده اندکی کوچک‌تر می‌شود. هدف از این کار، کاهش اثر پیکسل‌های مرزی، نواحی انسداد و خطاهای تطبیق در اطراف جسم است؛ زیرا مرزهای جسم معمولاً محل بروز خطاهای بیشتر در نقشه اختلاف مکانی هستند.

پس از تعیین ناحیه مورد نظر، پیکسل‌های نامعتبر یا دارای مقدار اختلاف مکانی نامناسب حذف می‌شوند. اگر تعداد پیکسل‌های معتبر از حداقل تعیین‌شده بیشتر باشد، یک مقدار نماینده برای اختلاف مکانی مانع محاسبه می‌شود. در حالت پیش‌فرض، از میانه مقادیر معتبر اختلاف مکانی استفاده می‌شود، زیرا میانه نسبت به نقاط پرت و خطاهای پراکنده مقاوم‌تر از میانگین ساده است.

با استفاده از رابطه هندسی دوربین استریو، عمق خام مانع از رابطه زیر به‌دست می‌آید:

\begin{equation}
	Z_k = \frac{fB}{d_k}
	\label{eq:app_stereo_depth}
\end{equation}

که در آن (f) فاصله کانونی بر حسب پیکسل، (B) فاصله مبنای دوربین و ($d_k$) اختلاف مکانی نماینده داخل ناحیه جسم است. مقدار ($Z_k$) در این رابطه، فاصله یا عمق مانع در راستای محور دید دوربین را نشان می‌دهد.

برای کاهش اثر مقادیر پرت، عمق خام هر مانع در یک بافر کوتاه‌مدت ذخیره می‌شود و پالایش مبتنی بر بازه بین‌چارکی \lr{IQR} روی مقادیر عمق انجام می‌گیرد. این پالایش روی عمق اعمال می‌شود، نه روی اختلاف مکانی، زیرا رابطه میان عمق و اختلاف مکانی غیرخطی است و پالایش مستقیم اختلاف مکانی می‌تواند باعث ایجاد سوگیری در تخمین فاصله شود. در این روش، ابتدا چارک اول و سوم مجموعه عمق‌های اخیر محاسبه می‌شوند:

\begin{equation}
	Q_1 = \mathrm{percentile}(Z,25),
	\qquad
	Q_3 = \mathrm{percentile}(Z,75)
	\label{eq:app_iqr_quartiles}
\end{equation}

سپس بازه بین‌چارکی برابر است با:

\begin{equation}
	\mathrm{IQR} = Q_3 - Q_1
	\label{eq:app_iqr_definition}
\end{equation}

نمونه‌هایی که خارج از بازه زیر قرار گیرند، به‌عنوان مقدار پرت در نظر گرفته می‌شوند:

\begin{equation}
	Q_1 - 1.5,\mathrm{IQR}
	\leq
	Z
	\leq
	Q_3 + 1.5,\mathrm{IQR}
	\label{eq:app_iqr_bounds}
\end{equation}

در صورت باقی‌ماندن نمونه معتبر، میانگین مقادیر باقی‌مانده به‌عنوان عمق پالایش‌شده آماری استفاده می‌شود. این مقدار سپس برای تشکیل مختصات دوبعدی مانع و به‌روزرسانی فیلتر کالمن به کار می‌رود.

زاویه نسبی مانع نسبت به محور دید دوربین از رابطه زیر محاسبه می‌شود:

\begin{equation}
	\theta_k =
	\tan^{-1}
	\left(
	\frac{u_{c,k}-c_x}{f_x}
	\right)
	\label{eq:app_bearing}
\end{equation}

که در آن ($u_{c,k}$) مختصات افقی مرکز جعبه محدودکننده، ($c_x$) مرکز اپتیکی تصویر و ($f_x$) فاصله کانونی افقی است. سپس مختصات تقریبی مانع در چارچوب دوربین به‌صورت زیر به‌دست می‌آید:

\begin{equation}
	x^{cam}_k = Z_k,
	\qquad
	y^{cam}_k = Z_k \tan(\theta_k)
	\label{eq:app_camera_coordinates}
\end{equation}

با در نظر گرفتن مکان نصب دوربین نسبت به لینک پایه، اندازه‌گیری نهایی برای فیلتر کالمن به‌صورت زیر ساخته می‌شود:

\begin{equation}
	x^{meas}_k = x^{cam}*k + x*{cam},
	\qquad
	y^{meas}*k = y*{cam} - y^{cam}_k.
	\label{eq:app_base_measurement}
\end{equation}

در پیاده‌سازی حاضر، ($x_{cam}=0.0~m$) و ($y_{cam}=0.0~m$) در نظر گرفته شده است. علامت منفی در مؤلفه (y) ناشی از تبدیل قرارداد محور تصویر و دوربین به قرارداد محور جانبی در چارچوب لینک پایه است.

%----------------------------------------------------------
\section{ماتریس‌های نویز}

در پیاده‌سازی انجام‌شده، نویز فرایند به‌صورت قطری تعریف شده است:

\begin{equation}
	\mathbf{Q} =
	\mathrm{diag}
	\left(
	q_{pos},;
	q_{vel},;
	q_{pos},;
	q_{vel}
	\right)
	\label{eq:app_kalman_Q}
\end{equation}

در تنظیمات مورد استفاده، مقدارهای این پارامترها برابر هستند با:

\begin{equation}
	q_{pos}=0.10,
	\qquad
	q_{vel}=0.80
	\label{eq:app_kalman_Q_values}
\end{equation}

مقدار بزرگ‌تر ($q_{vel}$) نسبت به ($q_{pos}$) به فیلتر اجازه می‌دهد تغییرات سرعت نسبی موانع را با انعطاف بیشتری دنبال کند. در مقابل، مقدار کوچک‌تر ($q_{pos}$) موجب می‌شود تغییرات موقعیت نسبت به نویزهای لحظه‌ای بیش از حد حساس نباشد.

نویز اندازه‌گیری نیز برای دو مؤلفه مکانی برابر فرض شده و به‌صورت زیر تعریف می‌شود:

\begin{equation}
	\mathbf{R} =
	r_{xy}\mathbf{I}_2
	\label{eq:app_kalman_R}
\end{equation}

در تنظیمات مورد استفاده:

\begin{equation}
	r_{xy}=0.10
	\label{eq:app_kalman_R_value}
\end{equation}

افزایش ($r_{xy}$) باعث می‌شود فیلتر به اندازه‌گیری‌های لحظه‌ای اعتماد کمتری داشته باشد و خروجی هموارتر شود. در مقابل، کاهش آن موجب واکنش سریع‌تر فیلتر به تغییرات اندازه‌گیری می‌شود. مقدار انتخاب‌شده در این پژوهش با هدف ایجاد تعادل میان نرمی خروجی و پاسخ‌گویی به حرکت نسبی موانع تنظیم شده است.

%----------------------------------------------------------
\section{مقداردهی اولیه}

هنگامی که یک مانع برای نخستین بار مشاهده می‌شود، فیلتر کالمن برای آن جسم ایجاد می‌گردد. مقدار اولیه موقعیت برابر با اندازه‌گیری جاری قرار داده می‌شود و سرعت اولیه صفر فرض می‌شود:

\begin{equation}
	\mathbf{x}_0 =
	\begin{bmatrix}
		x^{meas}_0 \
		0 \
		y^{meas}_0 \
		0
	\end{bmatrix}
	\label{eq:app_kalman_initial_state}
\end{equation}

همچنین کوواریانس اولیه حالت به‌صورت ماتریس همانی با ضریب (1.0) انتخاب می‌شود:

\begin{equation}
	\mathbf{P}_0 = \mathbf{I}_4
	\label{eq:app_kalman_initial_cov}
\end{equation}

این انتخاب نشان‌دهنده عدم قطعیت نسبتاً زیاد در زمان ایجاد مسیر جدید است. پس از چند مشاهده متوالی، فیلتر با دریافت اندازه‌گیری‌های جدید به‌تدریج مقدارهای موقعیت و سرعت را پایدارتر تخمین می‌زند.

%----------------------------------------------------------
\section{مرحله پیش‌بینی}

در هر گام، ابتدا ماتریس گذار حالت با مقدار ($\Delta t$) واقعی همان مانع به‌روزرسانی می‌شود. سپس مرحله پیش‌بینی فیلتر انجام می‌گیرد:

\begin{equation}
	\hat{\mathbf{x}}^{-}_k =
	\mathbf{F}*k \hat{\mathbf{x}}*{k-1}
	\label{eq:app_kalman_predict_state}
\end{equation}

\begin{equation}
	\mathbf{P}^{-}_k =
	\mathbf{F}*k \mathbf{P}*{k-1}\mathbf{F}_k^{T}
	+
	\mathbf{Q}
	\label{eq:app_kalman_predict_cov}
\end{equation}

در این رابطه، ($\hat{\mathbf{x}}^{-}_k$) و ($\mathbf{P}^{-}_k$) به‌ترتیب حالت و کوواریانس پیش‌بینی‌شده هستند. مرحله پیش‌بینی امکان می‌دهد حتی پیش از اعمال اندازه‌گیری جدید، وضعیت مانع بر اساس مدل حرکت سرعت ثابت تخمین زده شود.

%----------------------------------------------------------
\section{مرحله تصحیح}

پس از دریافت اندازه‌گیری جدید، نوآوری یا باقیمانده اندازه‌گیری به‌صورت زیر محاسبه می‌شود:

\begin{equation}
	\mathbf{y}_k =
	\mathbf{z}_k -
	\mathbf{H}\hat{\mathbf{x}}^{-}_k
	\label{eq:app_kalman_innovation}
\end{equation}

کوواریانس نوآوری برابر است با:

\begin{equation}
	\mathbf{S}_k =
	\mathbf{H}\mathbf{P}^{-}_k\mathbf{H}^{T}
	+
	\mathbf{R}
	\label{eq:app_kalman_innovation_cov}
\end{equation}

سپس بهره کالمن از رابطه زیر به‌دست می‌آید:

\begin{equation}
	\mathbf{K}_k =
	\mathbf{P}^{-}_k\mathbf{H}^{T}\mathbf{S}_k^{-1}
	\label{eq:app_kalman_gain}
\end{equation}

حالت و کوواریانس نهایی پس از تصحیح به‌صورت زیر محاسبه می‌شوند:

\begin{equation}
	\hat{\mathbf{x}}_k =
	\hat{\mathbf{x}}^{-}_k
	+
	\mathbf{K}_k\mathbf{y}_k
	\label{eq:app_kalman_correct_state}
\end{equation}

\begin{equation}
	\mathbf{P}_k =
	\left(
	\mathbf{I}
	-
	\mathbf{K}_k\mathbf{H}
	\right)
	\mathbf{P}^{-}_k
	\label{eq:app_kalman_correct_cov}
\end{equation}

در پیاده‌سازی عملی، این مراحل با استفاده از کلاس \lr{cv2.KalmanFilter} انجام می‌شوند. با این حال، منطق اصلی همان چرخه پیش‌بینی و تصحیح استاندارد فیلتر کالمن است.

%----------------------------------------------------------
\section{خروجی فیلتر}

پس از مرحله تصحیح، مؤلفه‌های حالت نهایی استخراج می‌شوند:

\begin{equation}
	x_k = \hat{\mathbf{x}}*k[0],
	\qquad
	v*{x,k} = \hat{\mathbf{x}}_k[1],
	\qquad
	y_k = \hat{\mathbf{x}}*k[2],
	\qquad
	v*{y,k} = \hat{\mathbf{x}}_k[3]
	\label{eq:app_kalman_output_components}
\end{equation}

سرعت کلی نسبی مانع نیز از رابطه زیر محاسبه می‌شود:

\begin{equation}
	v_k =
	\sqrt{
		v_{x,k}^{2}
		+
		v_{y,k}^{2}
	}
	\label{eq:app_kalman_speed}
\end{equation}

خروجی نهایی برای هر مانع شامل شناسه، کلاس، مختصات پالایش‌شده، فاصله پالایش‌شده، فاصله خام، زاویه نسبی، مؤلفه‌های سرعت نسبی، سرعت کلی و مقدار اختلاف مکانی خام است. فاصله خام مستقیماً از نقشه اختلاف مکانی به‌دست می‌آید، در حالی که فاصله پالایش‌شده از مؤلفه طولی حالت کالمن استخراج می‌شود.

لازم است توجه شود که سرعت‌های حاصل، سرعت‌های تخمینی در چارچوب لینک پایه و بر اساس تغییرات زمانی اندازه‌گیری‌های نسبی هستند؛ بنابراین نباید به‌عنوان سرعت مطلق جهانی مانع تفسیر شوند. برای مثال، اگر مانع در محیط ثابت باشد اما اسکوتر به سمت آن حرکت کند، فیلتر می‌تواند سرعت نسبی غیرصفر برای آن مانع تخمین بزند. این سرعت‌ها برای تصمیم‌گیری عامل یادگیری تقویتی، به‌ویژه برای تشخیص روند نزدیک‌شدن یا دورشدن موانع، مورد استفاده قرار می‌گیرند.

%----------------------------------------------------------
\section{مدیریت نبود اندازه‌گیری معتبر}

اگر در یک فریم، برای یک مانع مقدار اختلاف مکانی معتبر یا تعداد کافی پیکسل معتبر داخل ناحیه جعبه محدودکننده وجود نداشته باشد، آن مانع در خروجی همان فریم منتشر نمی‌شود. همچنین اگر عمق محاسبه‌شده خارج از بازه معتبر باشد، اندازه‌گیری رد می‌شود. در پیاده‌سازی انجام‌شده، بازه معتبر عمق برای موانع برابر (0.3~m) تا (6.0~m) در نظر گرفته شده است.

با وجود رد شدن اندازه‌گیری در یک فریم، وضعیت فیلتر مربوط به آن جسم تا مدت مشخصی در حافظه نگه‌داری می‌شود. اگر جسم در فریم‌های بعدی دوباره مشاهده شود، امکان ادامه پالایش زمانی وجود دارد. در پیاده‌سازی حاضر، اشیائی که برای مدت بیش از (5~s) مشاهده نشوند، از حافظه حذف می‌شوند.

این سازوکار از باقی‌ماندن طولانی‌مدت مسیرهای نامعتبر جلوگیری می‌کند و هم‌زمان اجازه می‌دهد وقفه‌های کوتاه‌مدت در آشکارسازی باعث حذف فوری وضعیت زمانی جسم نشود.

%----------------------------------------------------------
\section{جمع‌بندی}

الگوریتم ارائه‌شده در این پیوست، خروجی خام آشکارساز دوبعدی و نقشه اختلاف مکانی را به موقعیت و سرعت نسبی موانع تبدیل می‌کند. در این زنجیره، آشکارساز دوبعدی محل مانع را در تصویر مشخص می‌کند، نقشه اختلاف مکانی فاصله تقریبی آن را فراهم می‌سازد، پالایش \lr{IQR} اثر مقادیر پرت عمق را کاهش می‌دهد و فیلتر کالمن موقعیت دوبعدی و سرعت نسبی مانع را به‌صورت زمانی پایدارتر تخمین می‌زند.

این ساختار باعث می‌شود داده‌های مربوط به موانع برای استفاده در فضای مشاهده عامل تصمیم‌گیر و همچنین تحلیل وضعیت ایمنی مناسب‌تر شوند. با این حال، دقت خروجی نهایی همچنان به کیفیت تشخیص دوبعدی، کیفیت نقشه اختلاف مکانی، شرایط نوری، بافت سطح مانع و پایداری کالیبراسیون دوربین استریو وابسته است.
